\section{Threats to Validity}
\label{sec:threats}

\paragraph{Construct validity.}
\begin{itemize}[leftmargin=*]
\item \textbf{Seeded defects may not resemble real ones.} We mitigate this in three ways. The operators are derived from defect categories that published benchmark audits report~\cite{cuadron2025saber,zhu2025establishing}. Every mutant carries a kill witness, so none is equivalent. We also evaluate on real defects from the subject's history.
\item \textbf{The matching rules decide what counts as a detection.} The rules were fixed before any method ran, are the same for all methods, and are implemented once in shared scoring code.
\item \textbf{Scorability could look like cherry-picking.} The scorable set is computed from base runs only, before any mutant runs, and is identical for every method. Tables over all mutants are always reported alongside it.
\end{itemize}

\paragraph{Internal validity.}
\begin{itemize}[leftmargin=*]
\item \textbf{LLM errors are correlated across models.} This could make a shared misreading look like the truth. The anchored contrasts and the regime analysis address this, and the (A3) arm of the simulation measures the effect directly.
\item \textbf{Non-determinism.} We fix temperatures and seeds, pin model weight digests, fix the Ollama version, and cache every run with its digest.
\item \textbf{The Tell variants are generated by an LLM.} They paraphrase only the standing instruction and the task statement, and they were frozen once and are published.
\item \textbf{The hyper-parameters may be overfitted.} \specdx's thresholds were chosen by simulation calibrated to base runs only. The baselines' thresholds were chosen on the development split, which favours the baselines.
\item \textbf{The development gate is a second iteration on the same 20 development mutants.} An earlier version of the method was gated on them. The integrity rules of \cref{sec:eval-sim} limit the risk:
  \begin{itemize}
  \item every parameter of the current version was frozen before the gate;
  \item the development mutants are used once per version;
  \item the test split was never touched before the final study.
  \end{itemize}
\item \textbf{The population was relaxed after the screen.} The admission threshold on act rate was lowered from $0.6$ to $0.3$ during design, while the screen results were known. The screen reads no mutant and no truth, and the change was recorded in the pre-registration before any result of the current version.
\end{itemize}

\paragraph{External validity.}
\begin{itemize}[leftmargin=*]
\item \textbf{One subject.} The study covers one specification in one language and one domain.
\item \textbf{A narrow population.} It uses four local models of up to about 30B parameters from three families. Larger or hosted models might act more reliably and change the regime.
\item \textbf{Baselines on open models.} The LLM-based baselines B1 and B2 run on the same open models instead of the frontier models of their original papers, so their absolute performance may be lower than reported there. Each baseline's best model is selected on the development split, and the comparison at equal models is the fair one for our question.
\end{itemize}

\paragraph{Conclusion validity.}
\begin{itemize}[leftmargin=*]
\item \textbf{The test split is small.} The current operators yield 36 mutants, 16 of them in the test split. Two types have only two test mutants, so per-type McNemar tests have low power.
\item \textbf{Mitigations.} We report paired risk differences with bootstrap intervals and Holm-corrected exact tests, and we interpret effect sizes rather than $p$-values. Growing the operator sites is the most direct way to strengthen the conclusions, and we leave it to future work.
\end{itemize}
